\subsection{Продвинутая работа с камерой (только на дроне)}

Мы уже использовали камеру для поиска маркеров. Теперь разберемся, как работать с видеопотоком более гибко. Всё в этом разделе требует компьютера, подключённого к настоящему Пионеру.

\subsubsection*{Класс VideoStream}
Для получения <<сырого>> видеопотока (без автоматического декодирования в OpenCV формат, что иногда нужно для оптимизации) используется класс \texttt{VideoStream}.

\begin{codefile}[python]{python/camera_examples/camera_stream.py}
\end{codefile}

Этот код запускает передачу видео по Wi-Fi. Вы можете открыть этот поток в плеере VLC (адрес обычно \texttt{rtsp://192.168.4.1:8554/front}). В симуляторе \texttt{VideoStream} --- только обёртка над той же камерой: поток наружу браузер не отдаёт, открыть его в VLC не получится.

\subsubsection*{Захват кадров и фотографирование}
Чтобы сделать фото, нужно получить текущий кадр и сохранить его в файл с помощью \texttt{cv2.imwrite(имя\_файла, кадр)}. В симуляторе \texttt{imwrite} возвращает \texttt{True}, но файла не создаёт: файловая система Python в браузере живёт в памяти вкладки.

\subsubsection*{Калибровка камеры}
Для точных измерений (например, расстояния до маркера через \texttt{solvePnP}, раздел~\ref{sec:python-aruco-pose}) нужно знать оптические параметры камеры: матрицу камеры и коэффициенты дисторсии. Камера симулятора идеальная, её матрицу программе отдаёт \texttt{cv2.FileStorage}; калибровка нужна только настоящей камере, и подробно она разобрана в главе~\ref{sec:python-map-creation}.

Процесс калибровки:
\begin{enumerate}
    \item Распечатать шахматную доску (Chessboard).
    \item Сделать серию снимков доски под разными углами.
    \item Запустить скрипт калибровки (\texttt{camera\_examples/camera\_calibration.py}; его разбор --- в главе~\ref{sec:python-map-creation}). Он вычислит параметры и сохранит их в файл \texttt{data.yml}: матрица камеры --- узел \texttt{mtx}, коэффициенты дисторсии --- узел \texttt{dist}. Этот же файл читают примеры с оценкой позы из предыдущих разделов.
\end{enumerate}

\begin{infobox}[Зачем это нужно?]
Без калибровки прямые линии на изображении могут казаться изогнутыми (эффект <<рыбьего глаза>>), что приведет к ошибкам в навигации.
\end{infobox}

\begin{taskbox}[Фотограф]
Только на дроне. Напишите программу, которая показывает видео с дрона в окне \texttt{cv2.imshow}. При нажатии клавиши \texttt{S} текущий кадр сохраняется в файл \texttt{photo\_N.jpg}, где N --- порядковый номер фото, начиная с 1; клавиша \texttt{Q} завершает программу и закрывает окно. Пустой кадр (\texttt{None}) не сохранять.
\expected{после трёх нажатий \texttt{S} рядом с программой лежат \texttt{photo\_1.jpg}--\texttt{photo\_3.jpg}, в терминале --- три строки «Сохранено: photo\_N.jpg»; после \texttt{Q} окно закрывается без ошибки.}
\end{taskbox}
